\section{Introduction}
\label{sec:introduction}

Least-squares methods provide a systematic way to obtain an approximate
solution when a linear system is overdetermined or when a set of
observations cannot be represented exactly by a prescribed model. Their
reliability depends not only on the mathematical formulation but also on the
numerical properties of the algorithms used to compute the solution
\cite{golub2013matrix,trefethen1997numerical}.

This work investigates two least-squares problems specified in the
Computational Work 3 assignment. First, a linear system
$A\mathbf{x}\approx\mathbf{b}$, with $A\in\mathbb{R}^{5\times4}$, is addressed
through the reduced QR factorization of $A$. Second, six data points are
fitted with a polynomial of degree at most two, expressed as a linear
least-squares problem involving a design matrix $X\in\mathbb{R}^{6\times3}$.
In both cases, the data are constructed according to the assignment
instructions from the digits of the student's enrollment number.

The computational objective is to connect these formulations to their
implementation. The main numerical routines—including inner products,
vector and matrix operations, the Modified Gram--Schmidt QR factorization,
backward substitution, and error indicators—are implemented explicitly in
Python. NumPy is retained for the array construction required by the
assignment and for generating plotting coordinates, while Matplotlib is used
to visualize the polynomial fit. No high-level NumPy or SciPy solver is used
to calculate the QR factorization or the least-squares coefficients. The
implementation and selected code excerpts are provided in the appendix.

The computed solutions are assessed through the QR reconstruction error, the
orthogonality error of $Q$, the least-squares optimality condition, and, for
the polynomial fit, the mean squared error. These measures serve different
purposes: the first indicators assess numerical consistency of the
algorithms, whereas the mean squared error quantifies the discrepancy between
the observations and the fitted model. The analysis also examines the
conditions under which each least-squares solution is unique, while
acknowledging the numerical limitations of Modified Gram--Schmidt without
column pivoting \cite{higham2002accuracy}.
